\section{Have you ever drunk the Silence?}

\begin{quotex}
Concentration without effort ... is your life tossed to and fro by random events, thoughts, feelings? Or do you live life consciously? It begins with Silence ... 
\end{quotex}

\textbf{Rene Guenon} claimed that at times when the authorities had lost the inner meaning of things, initiates would pose as jugglers or horse traders. That way, they could travel from village to village, under cover as it were, to meet with other initiates. One can imagine them carrying Tarot cards as a teaching tool, since they appear to be a harmless game, and are much more compact than transporting a library. That is how I see the first card, Le Bateleur (the Juggler or Magician), as an itinerant initiate. The name of the card is a French pun on ``the low deceives you'' (``le bas te leurre''), but the initiates are not deceived.

\textbf{Valentin Tomberg} relates this card to ``concentration without effort'', reminiscent of Taoism, which is the necessary first step on the journey through the deck. Our Friend writes:

\begin{quotex}
Concentration \emph{without effort}, which means there is nothing to suppress and where contemplation becomes as natural as breathing and the beating of the heart, is the state of consciousness --- of the intellect, the imagination, the feelings, and the will --- a state of perfect calm, accompanied by the complete relaxation of the nerves and muscles of the body. It is the deep silence of desires, concerns, imagination, memory, and discursive thought. We would say that the entire being has become like the surface of calm waters reflecting the immense presence of the starry sky and its inexpressible harmony. And the waters are deep, oh how deep! And the silence increases, always increasing, what SILENCE! Its growth takes place in regular waves which pass, one after the other, through your being: one wave of silence followed by another wave of deeper silence, then yet another wave of even deeper silence ... Have you ever \emph{drunk the silence}. If so, you know what concentration without effort it.
\end{quotex}



\flrightit{Posted on 2010-10-06 by Cologero}

